% !TeX root = ../sections/02_exercises.tex

\begin{exercise}
	R-10. 環 $\mathbb{Z}[X,Y]$ のイデアル
	\[
	I=(X^2-2,\;Y^2-3,\;XY)
	\]
	は $X^2-2,\;Y^2-3,\;XY$ で生成されたイデアルとする。
	\[
	A=\mathbb{Z}[X,Y]/I
	\]
	と置き，
	\[
	\overline{X}=X+I\in A,
	\qquad
	\overline{Y}=Y+I\in A
	\]
	と表す。
	\begin{enumerate}
		\item $\mathbb{Z}$-加群として
		\[
		A\cong \mathbb{Z}\oplus \mathbb{Z}/2\mathbb{Z}\oplus \mathbb{Z}/3\mathbb{Z}
		\]
		であることを示せ。
		\item $B=(3\overline{X}+2\overline{Y})$ を
		$3\overline{X}+2\overline{Y}$ で生成された $A$ の単項イデアルとするとき，
		\[
		B=6\mathbb{Z}\oplus\mathbb{Z}(3\overline{X}+2\overline{Y})
		\]
		を確かめよ。
		また，$\mathbb{Z}$-加群として $A/B$ の不変系を求めよ。
		\item $C=(3\overline{X},2\overline{Y})$ を
		$3\overline{X}$ と $2\overline{Y}$ で生成された $A$ のイデアルとするとき，
		$\mathbb{Z}$-加群として $A/C$ の不変系を求めよ。
	\end{enumerate}
\end{exercise}

\begin{background}
	商環
	\[
	A=\mathbb{Z}[X,Y]/I
	\]
	では，
	\[
	\overline{X}^2=2,
	\qquad
	\overline{Y}^2=3,
	\qquad
	\overline{X}\overline{Y}=0
	\]
	という関係が成り立つ。
	
	したがって，$A$ の元はまず
	\[
	a+b\overline{X}+c\overline{Y}
	\]
	の形に整理することを考える。
	
	ただし，この設問は注意が必要である。
	与えられた関係式を環の等式として厳密に使うと，
	\[
	6=\overline{X}^2\overline{Y}^2
	=(\overline{X}\overline{Y})^2=0
	\]
	のような追加関係も現れる。
	そのため，問題文の意図としてどの $\mathbb{Z}$-加群表示を用いるかを確認しながら進める必要がある。
	
	不変系を求める基本方針は，$A$ を $\mathbb{Z}$-加群として表示し，
	その中で $B$ や $C$ を部分加群として書き，生成元行列の Smith 標準形を求めることである。
\end{background}

\begin{strategy}
	まず，商環 $A$ の関係式
	\[
	\overline{X}^2=2,
	\qquad
	\overline{Y}^2=3,
	\qquad
	\overline{X}\overline{Y}=0
	\]
	を用いて，任意の元を
	\[
	a+b\overline{X}+c\overline{Y}
	\]
	の形に整理する。
	
	次に，$B=(3\overline{X}+2\overline{Y})$ については，
	生成元
	\[
	g=3\overline{X}+2\overline{Y}
	\]
	に対して，
	\[
	g\cdot 1,\qquad
	g\cdot\overline{X},\qquad
	g\cdot\overline{Y}
	\]
	を計算する。
	これにより，$B$ が $\mathbb{Z}$-加群としてどの元で生成されるかを調べる。
	
	同様に，$C=(3\overline{X},2\overline{Y})$ については，
	\[
	3\overline{X},\qquad 2\overline{Y}
	\]
	およびそれらに $\overline{X},\overline{Y}$ を掛けた元を調べる。
	
	最後に，$A$ の $\mathbb{Z}$-加群表示に対して，
	$B$ や $C$ の生成元を列ベクトルで表し，Smith 標準形を求める。
	その対角成分から
	\[
	A/B,\qquad A/C
	\]
	の不変系を読み取る。
\end{strategy}